\scanpage{56}
\margin{55}
- \mspage{75} -
\hfill n°269

a $y\in S_1\cap F_2=S_1\cap F_1\cap F_2=S_1\cap F=\{0\}$. On a donc $S_1\cap S_2^{\mathrm o}=\{0\}$, et de même $S_2\cap S_1^{\mathrm o}=\{0\}$. Comme l'un des deux sous-espaces $F_1,F_2$, par exemple $F_1$, est de dimension finie, $S_1$ est de dimension finie, donc $S_1^{\mathrm o}$ est de codimension finie, et par conséquent $S_2$ est de dimension finie puisque $S_2\cap S_1^{\mathrm o}=\{0\}$; de plus ceci montre que l'on a $\dim S_2\leq\operatorname{codim}S_1^{\mathrm o}=\dim S_1$ ; de même $\dim S_1\leq\dim S_2$, d'où $\dim S_1=\dim S_2$. La prop. 2,a) montre alors que $S=S_1+S_2$ est non isotrope.

Ceci étant, l'orthogonal $N$ de $S$ est non isotrope (cor. de la prop.1) et contient $F$. Le cor. à la prop.2 montre donc qu'il existe un sous espace $G$ totalement isotrope (resp. totalement singulier) de $N$ tel que $\dim G=\dim F$, que $G\cap F=\{0\}$ et que $G+F$ soit non dégénéré. Ainsi d) est vérifiée par $G$. On satisfera alors à a) et b) en prenant pour $H$ l'orthogonal de $G+F$ dans $N$. Quant à c), l'on remarque que, comme $H$ est orthogonal à $F_1=S_1+F$, il n'y a aucun vecteur isotrope (resp. singulier) non nul dans $H$ en vertu de ce qui a été remarqué au début de la démonstration et du fait que $H\cap F_1=\{0\}$. Enfin les assertions relatives aux dimensions ont été démontrées en cours de route, ou s'en déduisent trivialement.

COROLLAIRE 1 - \uline{Les hypothèses étant celles de la prop.3, deux sous-espaces totalement isotropes} (resp. \uline{totalement singuliers}) \uline{maximaux de dimensions finies ont même dimension. Pour tout sous-espace totalement isotrope} (resp. \uline{totalement singulier}) \uline{maximal} $F$ \uline{de dimension finie, il en existe un autre} $F'$ \uline{tel que} $F\cap F'=\{0\}$, \uline{et dans ces conditions} $F+F'$ \uline{est non isotrope.}
